% ==================================================================
% id:       comb-019
% topic:    קומבינטוריקה
% subtopic: עקרון הכפל וסידורים
% solution: full
% ==================================================================
על מדף מסדרים 9 ספרים: 4 כחולים, 3 אדומים ו-2 צהובים. בכמה דרכים אפשר לסדר אותם:
\begin{enumerate}
\item כאשר אפשר להבחין בין כל תשעת הספרים?
\begin{solbox}
סידור של 9 עצמים שונים בשורה: $9!=\ans{362{,}880}$.
\end{solbox}
\item כאשר אי אפשר להבחין בין ספרים מאותו צבע?
\hint{כמה סידורים של ספרים שונים נראים בדיוק אותו דבר, כשמסתכלים רק על הצבעים?}
\begin{solbox}
עכשיו סידור נקבע רק לפי רצף הצבעים, למשל $\text{\LR{BBRYBRBYR}}$. כל רצף צבעים כזה מתאים לבדיוק $4!\cdot3!\cdot2!$ סידורים של ספרים שונים: אפשר לסדר את 4 הכחולים ב-4 המקומות הכחולים ב-$4!$ דרכים, את האדומים ב-$3!$ ואת הצהובים ב-$2!$. לכן
\[
\frac{9!}{4!\,3!\,2!}=\ans{1260}.
\]
(זה אותו רעיון כמו במקדם הבינומי: סופרים כאילו הכול שונה, ומחלקים במספר הפעמים שכל אפשרות נספרה.)
\end{solbox}
\end{enumerate}
